\documentclass[10pt,a4paper]{amsart}
\usepackage[margin=19mm]{geometry}
\usepackage{amsmath,amssymb,mathtools}
\usepackage[scr=rsfs,cal=euler]{mathalfa}
\usepackage{xfrac}
\usepackage[unicode,hidelinks,bookmarks=false]{hyperref}
\newcommand{\ie}{i.e.\ }
\newcommand{\cf}{cf.\ }
\newcommand{\resp}{respectively\ }
\newcommand{\Subsection}{\subsection}
\providecommand{\nobreakdash}{\nobreak\hbox{-}}
\setlength{\emergencystretch}{2em}
\multlinegap0pt
\usepackage{amssymb}
\usepackage{mathtools}
\newtheorem{proposition}{Proposition}
\title{The second minimum weight of Grassmann codes}
\author{Mrinmoy Datta, Tiasa Dutta}
\date{Source: arXiv:2508.18844v2 (CC BY 4.0)}
\begin{document}
\maketitle

\noindent\emph{Source and licence.} The second minimum weight of Grassmann codes. Version arXiv:2508.18844v2 (CC BY 4.0), distributed by arXiv under the Creative Commons Attribution 4.0 International licence, http://creativecommons.org/licenses/by/4.0/, as recorded in the arXiv licence metadata for this exact version and shown on the abstract page https://arxiv.org/abs/2508.18844v2. Full licence notice: \texttt{licenses/arxiv-2508.18844-CC-BY-4.0.txt}.\par\smallskip

\noindent\textit{Editorial scope.} Counted unit: the article's proposition ineq (source0.tex lines 644--652). The article's complete original proof of it is delivered with it (source0.tex lines 654--672, one \texttt{proof} environment of 19 lines). No mathematics is written, repaired or re-derived: the statement and the proof are the article's own lines, in the article's own notation, and no retained line is substituted or re-flowed.  Retained uncounted as the source-local context the proof needs: lines 472--479.  Nothing was omitted from any retained span, so the package carries no omission record.  No retained line contains a citation, so the wrapper carries no bibliography and no bibliography entry.  No retained line contains a figure environment, an image-inclusion command or drawing code, so no image is delivered and the package declares no asset dependency.  Editorial formatting only, in addition: an AMS \texttt{amsart} preamble that restates the article's own declarations and macros used by the retained lines, namely the environments \texttt{proposition} on separate counters, together with the standard packages those lines need.  The retained environments print in this document as Proposition 1 in order of appearance, whereas the article prints the same items with the numbering of its own full document.  The complete native source of the article as distributed in the arXiv e-print for this exact version is bundled unchanged as \texttt{sources/183903/source0.tex}.
\begin{proposition}\label{identities}
    For positive integers $\ell, m$ satisfying $1 \le \ell \le m$, we have,
    \begin{enumerate}
        \item[(a)] ${m \brack \ell}_q = {m \brack m - \ell}_q$. 
        \item[(b)] ${m \brack \ell}_q = {m - 1 \brack \ell}_q + q^{m- \ell}{m - 1 \brack \ell - 1}_q$.
        \item[(c)] ${m \brack \ell}_q = \frac{q^m - 1}{q^{m - \ell} - 1} {m - 1 \brack \ell}_q$
    \end{enumerate}
\end{proposition}

\begin{proposition}\label{ineq}
    Let $\ell, m$ be positive integers with $1 \le \ell \le m-1$. We have 
    \begin{enumerate}
        \item[(a)] $e (\ell, m-1) \left(\frac{q^m -1}{q^{m-\ell} - 1}\right) < e (\ell, m).$
        \item[(b)] ${m - 1 \brack \ell}_q + q^{m-\ell} e (\ell -1, m-1) = e (\ell, m). $
        \item[(c)] If $2 \le \ell \le m-2$, then $e' (\ell, m-1) \left(\frac{q^m -1}{q^{m-\ell} - 1}\right) < e' (\ell, m).$
        \item[(d)] If $2 \le \ell \le m-2$, then ${m - 1 \brack \ell}_q + q^{m-\ell} e' (\ell -1, m-1) = e' (\ell, m)$.
    \end{enumerate}
\end{proposition}

\begin{proof}
    \
    \begin{enumerate}
        \item[(a)] Note that \begin{align*}
        e (\ell, m-1) \left(\frac{q^m -1}{q^{m-\ell} - 1}\right) &=\left(\frac{q^m - 1}{q^{m-\ell} - 1} \right) \left( {m - 1 \brack \ell}_q - q^{\ell(m-1 - \ell)} \right) \\
        &= {m \brack \ell}_q - \left(\frac{q^m - 1}{q^{m-\ell} - 1}\right) q^{\ell(m-1 - \ell)} \\
        &< {m \brack \ell}_q - q^{\ell (m-\ell)} = e(\ell, m).
    \end{align*}
    The equality above follows from Proposition \ref{identities} (c). 
    while the inequality on the last line is equivalent to the inequality $(q^m - 1) > q^\ell (q^{m - \ell} - 1) = q^m - q^\ell$, which is trivially true as $\ell \ge 1$ and $q$ is a prime power. 
    \item[(b)] The left-hand side of the assertion is given by
        \begin{equation*}
        {m - 1 \brack \ell}_q + q^{m - \ell} \left({ m - 1 \brack \ell - 1}_q - q^{(\ell-1)(m - \ell)} \right) \\
       =  {m \brack \ell}_q - q^{\ell(m- \ell)} = e(\ell, m). 
    \end{equation*} 
    The first equality follows from  Proposition \ref{identities} (b). 
    \end{enumerate}
    We leave the proof of parts (c) and (d), which can be easily deduced in a way similar to parts (a) and (b) respectively, to the reader. 
\end{proof}

\end{document}
